\subsubsection{跨性别者就医流程实务：从挂号到复诊的每一步}

上文介绍了性别肯定医疗的构成（心理支持、社会过渡、激素、手术）与中国的指南现状。本节补上最常被忽略的一层：\textbf{具体去医院的流程是什么样、每一步可能遇到什么、怎么减少摩擦}。

\textbf{一、初诊与评估}
\begin{itemize}
\item \textbf{精神科/心理科评估}：国内的性别肯定医疗路径通常从精神科或心理科的评估开始，出具与"性别不一致"相关的评估意见。这一步的现实状况是——并非所有医院都开展该项业务，需要提前确认接诊机构；评估本身可能不止一次面诊；
\item \textbf{准备材料}：通常需要身份证明；部分机构要求家长知情或陪同（成年人与未成年人的要求不同，务必先电话确认）。把要求问清楚再跑一趟，能省下大量往返；
\item \textbf{如何问}：致电医院总机或心理科时，可以直接问"贵院是否可以做性别不一致的评估与转介"。直接询问比绕圈子更有效——这类信息本身不是隐私问题，也不构成任何"记录"。
\end{itemize}

\textbf{二、激素治疗}
\begin{itemize}
\item \textbf{正规渠道优先}：由内分泌科或相关专科医生开具处方、定期复查肝肾功能、血脂、血细胞比容与激素水平，是激素治疗的基本安全框架；
\item \textbf{自行用药的风险要说清楚}：非正规渠道获取的激素存在剂量不可控、成分不明、缺乏监测三重风险。常见后果包括肝功能损伤、血栓风险升高（尤其雌激素合用吸烟）、骨密度变化与电解质紊乱。这些风险不是用来吓阻——而是说明\textbf{定期抽血检查的价值}：它能把"盲用"变成"可控"；
\item \textbf{主动告知用药史}：无论激素来自何处，就医时都应如实告知在用药物与剂量。隐瞒导致的药物相互作用（如与抗凝药、降压药、降糖药）风险远大于"被批评"的尴尬。
\end{itemize}

\textbf{三、手术与住院}
\begin{itemize}
\item 手术通常有明确的前置条件（如年龄、社会过渡时长、精神科评估意见），流程较长且需要提前规划时间。
\item \textbf{住院期间的实务}：入院登记、病历、床头卡与日常交班都可能涉及称呼问题。可以主动向主管护士说明希望被如何称呼、是否使用特定代词；不是所有医院都有统一规范，但\emph{主动沟通通常比沉默更能得到配合}。
\item \textbf{隐私与陪伴}：涉及隐私部位的手术与检查，可要求由指定性别的工作人员陪同或操作（在人力允许范围内）；如希望伴侣陪同，可在入院时说明与协商。
\item \textbf{术后随访}：手术不是终点，术后功能恢复、复查与心理调适都需要持续安排（本书《生殖器官的个体差异与跨性别视角》一节也谈到身体体验层面的议题）。
\end{itemize}

\textbf{四、日常就医中的三个小建议}
\begin{itemize}
\item \textbf{不必要的信息不主动提供，必要的需求要说清楚}：比如"我最近在服用激素，需要复查肝功能和激素水平"——只说与诊疗相关的部分，其他细节可自行决定是否披露；
\item \textbf{找得到友善的医生就固定下来}：一位熟悉你的医生能显著降低每次沟通的成本；
\item \textbf{遇到不友善对待怎么办}：医院有投诉渠道，卫健行政部门也有受理途径；感到被歧视时的记录（时间、科室、具体言行）会让后续处理更容易。要不要投诉由本人决定，但\emph{知道路径存在}本身就是一种保护。
\end{itemize}

\noindent\textit{【编者注】}本节只讲流程与实务，不做制度评价。跨性别者就医的摩擦往往来自"流程没有为他准备答案"，而多数摩擦可以通过提前确认信息、主动说明需求来减少——这也是本节的目的。
